\chapter{AES-CMAC from scratch}
\label{ch:cmac}

\chapterbrief
  {Build the message authentication code DESFire EV2 runs on.}
  {Immutable objects and constructor invariants.}
  {The RFC 4493 known-answer vectors passing.}

\section{Read this before you write any cryptography}

Everything from here on is \textbf{verified against known-answer test vectors
before it touches hardware}. This is not caution for its own sake.

Cryptographic code has a property no other code you have written has:
\textbf{wrong output is indistinguishable from right output.} A broken CRC gives
you a mismatch you can see. A broken CMAC gives you sixteen bytes of
perfect-looking randomness --- exactly like a correct CMAC. You cannot eyeball
it, you cannot reason about it from a failed transaction, and if you build three
more layers before testing you will have no idea which one is lying.

On real hardware it is worse. A bad MAC does not merely fail: the card may
increment an internal error counter, and enough consecutive failures can
permanently lock a key.

\begin{center}\itshape\color{brewerred}
Vectors first. Silicon second.
\end{center}

\section{What a MAC is}

A \textbf{Message Authentication Code} answers one question: \emph{was this
message written by someone holding the secret key, and has it been altered?}

It is a keyed checksum. A CRC tells you a message survived the wire; a MAC tells
you it survived an \emph{adversary}. Without the key you cannot compute a valid
MAC, so you cannot forge or tamper.

\begin{analogy}{The wax seal}
A wax seal on a letter. Anyone can read the letter, but only the person with the
signet ring can produce the seal, and breaking it to alter the contents is
obvious.

The ring is the key; the seal is the MAC. Crucially the seal does not \emph{hide}
the letter --- MACs authenticate, they do not encrypt. Those are separate jobs,
and DESFire uses both.
\end{analogy}

\textbf{CMAC} (RFC 4493 / NIST SP 800-38B) builds a MAC out of a block cipher
you already have. Roughly: it is CBC-mode encryption where you throw away
everything except the last block. That final block depends on every bit of the
message and on the key, so it works as a fingerprint.

The subtlety CMAC adds --- and the reason it is not simply ``CBC-MAC'' --- is how
it handles the \textbf{last block}, where naive CBC-MAC has a forgery weakness.
CMAC derives two subkeys, $K_1$ and $K_2$, and XORs one of them into the final
block depending on whether the message needed padding. That is the whole trick.

\section{The algorithm}

\subsection{Part 1 --- subkey generation (RFC 4493 §2.3)}

\begin{trace}
L  = AES-ECB(K, 0^128)                 encrypt an all-zero block
K1 = L << 1;   if MSB(L)  was set:  K1 ^= Rb
K2 = K1 << 1;  if MSB(K1) was set:  K2 ^= Rb
\end{trace}

\begin{itemize}
  \item $R_b = $ \hex{0x87}, the constant from AES's finite-field polynomial.
        ``\code{\^{}= Rb}'' means XOR \emph{the last byte} with \hex{0x87}.
  \item ``MSB set'' means the top bit of the \emph{first} byte:
        \code{(x[0] \& 0x80) != 0}.
  \item \code{<< 1} shifts the whole 128-bit block left one bit --- the bit that
        falls off byte $i{+}1$'s top becomes byte $i$'s bottom.
\end{itemize}

That last one is the only fiddly part, and in C you would write it identically:

\begin{jcode}
byte[] out = new byte[16];
int carry = 0;
for (int i = 15; i >= 0; i--) {      // walk BACKWARDS, low byte first
    int b = in[i] & 0xFF;
    out[i] = (byte) ((b << 1) | carry);
    carry = (b >> 7) & 1;            // the bit shifted out of this byte
}
\end{jcode}

You walk from the last byte to the first because a carry propagates
\emph{upward}, toward the more significant end.

\subsection{Part 2 --- the MAC (RFC 4493 §2.4)}

\begin{center}
\begin{tikzpicture}[font=\scriptsize\sffamily, node distance=5mm,
  blk/.style={draw=brewergrey, rounded corners=1pt, minimum width=13mm,
              minimum height=6mm, inner sep=2pt},
  op/.style={draw=brewernavy, circle, inner sep=1pt, minimum size=4.5mm},
  aes/.style={draw=brewernavy, fill=brewernavy!8, rounded corners=1pt,
              minimum width=13mm, minimum height=6mm}]
  \node[blk] (m1) {$M_1$};
  \node[blk, right=14mm of m1] (m2) {$M_2$};
  \node[blk, right=14mm of m2, draw=breweramber] (ml) {$M_{last}$};

  \node[op, below=5mm of m1] (x1) {$\oplus$};
  \node[op, below=5mm of m2] (x2) {$\oplus$};
  \node[op, below=5mm of ml] (xl) {$\oplus$};

  \node[aes, below=4mm of x1] (a1) {AES$_K$};
  \node[aes, below=4mm of x2] (a2) {AES$_K$};
  \node[aes, below=4mm of xl] (al) {AES$_K$};

  \node[left=7mm of x1, font=\scriptsize] (zero) {$0^{128}$};

  \draw[-{Stealth[length=3pt]}] (m1) -- (x1);
  \draw[-{Stealth[length=3pt]}] (m2) -- (x2);
  \draw[-{Stealth[length=3pt]}] (ml) -- (xl);
  \draw[-{Stealth[length=3pt]}] (zero) -- (x1);
  \draw[-{Stealth[length=3pt]}] (x1) -- (a1);
  \draw[-{Stealth[length=3pt]}] (x2) -- (a2);
  \draw[-{Stealth[length=3pt]}] (xl) -- (al);
  \draw[-{Stealth[length=3pt]}] (a1) -- (x2);
  \draw[-{Stealth[length=3pt]}] (a2) -- (xl);
  \node[below=4mm of al, font=\scriptsize\bfseries, color=brewergreen] (t) {CMAC};
  \draw[-{Stealth[length=3pt]}] (al) -- (t);

  \node[above=4mm of ml, align=center, font=\tiny, color=breweramber!80!black]
       {full block: $\oplus K_1$\\ padded: $\oplus K_2$};
\end{tikzpicture}
\end{center}

\begin{trace}
split the message into 16-byte blocks M_1 ... M_n
the LAST block is special:
    if it is exactly 16 bytes:   M_last = M_n XOR K1
    otherwise (incl. empty):     pad it (append 0x80 then 0x00s to 16)
                                 M_last = padded XOR K2
then CBC-MAC:
    X = 0^128
    for each block except the last:   X = AES(K, X XOR M_i)
    T = AES(K, X XOR M_last)          <- the CMAC
\end{trace}

The $K_1$ versus $K_2$ choice is what makes ``a 16-byte message'' and ``a 15-byte
message padded to 16'' produce different MACs. Without it, an attacker could
exploit the ambiguity.

\section{Build it}

First, thin AES helpers. This is the one piece handed to you complete ---
reimplementing AES itself teaches you nothing here, and the JDK's is correct and
fast.

\begin{jcode}
package com.midnightbrewer.hardware.crypto;

import javax.crypto.Cipher;
import javax.crypto.spec.IvParameterSpec;
import javax.crypto.spec.SecretKeySpec;

/**
 * Thin AES-128 helpers over the JDK's crypto provider.
 *
 * NoPadding everywhere: DESFire does its own ISO 7816-4 padding, and letting
 * the JDK add PKCS#5 on top would corrupt every frame.
 */
public final class Aes {

    public static final int BLOCK = 16;

    private Aes() { }

    public static byte[] ecbEncrypt(byte[] key, byte[] block) {
        try {
            Cipher c = Cipher.getInstance("AES/ECB/NoPadding");
            c.init(Cipher.ENCRYPT_MODE, new SecretKeySpec(key, "AES"));
            return c.doFinal(block);
        } catch (Exception e) {
            throw new IllegalStateException("AES-ECB failed", e);
        }
    }

    public static byte[] cbcEncrypt(byte[] key, byte[] iv, byte[] data) {
        try {
            Cipher c = Cipher.getInstance("AES/CBC/NoPadding");
            c.init(Cipher.ENCRYPT_MODE, new SecretKeySpec(key, "AES"),
                   new IvParameterSpec(iv));
            return c.doFinal(data);
        } catch (Exception e) {
            throw new IllegalStateException("AES-CBC encrypt failed", e);
        }
    }

    public static byte[] cbcDecrypt(byte[] key, byte[] iv, byte[] data) {
        try {
            Cipher c = Cipher.getInstance("AES/CBC/NoPadding");
            c.init(Cipher.DECRYPT_MODE, new SecretKeySpec(key, "AES"),
                   new IvParameterSpec(iv));
            return c.doFinal(data);
        } catch (Exception e) {
            throw new IllegalStateException("AES-CBC decrypt failed", e);
        }
    }

    public static byte[] xor(byte[] a, byte[] b) {
        byte[] out = new byte[a.length];
        for (int i = 0; i < a.length; i++) {
            out[i] = (byte) (a[i] ^ b[i]);
        }
        return out;
    }
}
\end{jcode}

Now the part you write.

\begin{jcode}
package com.midnightbrewer.hardware.crypto;

/**
 * AES-CMAC (RFC 4493) -- the MAC DESFire EV2 uses. The JDK has no CMAC, so
 * this is the from-scratch part; you are not reimplementing AES, you are
 * building the CMAC construction on top of Aes.ecbEncrypt.
 *
 * OOP on show: the key and the two derived subkeys are computed ONCE in the
 * constructor and kept private. compute() never mutates them. That makes this
 * an immutable object with an invariant established at birth: "if this object
 * exists, its subkeys are correct and match its key."
 */
public final class AesCmac {

    private static final int BLOCK = 16;
    private static final byte RB = (byte) 0x87;  // AES field polynomial

    private final byte[] key;
    private final byte[] k1;
    private final byte[] k2;

    public AesCmac(byte[] key) {
        if (key.length != BLOCK) {
            throw new IllegalArgumentException("AES-128 key must be 16 bytes");
        }
        this.key = key.clone();
        byte[][] subkeys = generateSubkeys(key);
        this.k1 = subkeys[0];
        this.k2 = subkeys[1];
    }

    // ---- RFC 4493 section 2.3 ----
    private static byte[][] generateSubkeys(byte[] key) {
        // TODO 1: byte[] l = Aes.ecbEncrypt(key, new byte[BLOCK]);
        // TODO 2: byte[] k1 = leftShift(l);
        //         if ((l[0] & 0x80) != 0) k1[15] ^= RB;
        // TODO 3: byte[] k2 = leftShift(k1);
        //         if ((k1[0] & 0x80) != 0) k2[15] ^= RB;
        // TODO 4: return new byte[][]{ k1, k2 };
        return new byte[][]{ new byte[BLOCK], new byte[BLOCK] };
    }

    /** Shift a 16-byte block left by ONE bit, as a single 128-bit number. */
    private static byte[] leftShift(byte[] in) {
        // TODO: walk from the LAST byte to the first, carrying the top bit
        //       of each byte down into the next. (See the chapter text.)
        return new byte[BLOCK];
    }

    // ---- RFC 4493 section 2.4 ----
    public byte[] compute(byte[] message) {
        // TODO 1: how many blocks, and is the last one full?
        //   if (message.length == 0) { n = 1; lastComplete = false; }
        //   else { n = ceil(message.length / 16.0);
        //          lastComplete = (message.length % 16 == 0); }

        // TODO 2: build M_last.
        //   if lastComplete: mLast = Aes.xor(lastBlock, k1)
        //   else: pad into a 16-byte buffer -- copy the leftover bytes,
        //         set padded[leftoverLen] = (byte) 0x80, rest stays zero;
        //         mLast = Aes.xor(padded, k2)

        // TODO 3: CBC-MAC through the blocks:
        //   byte[] x = new byte[BLOCK];
        //   for (int i = 0; i < n - 1; i++)
        //       x = Aes.ecbEncrypt(key, Aes.xor(x, i-th 16-byte slice));
        //   return Aes.ecbEncrypt(key, Aes.xor(x, mLast));
        return new byte[BLOCK];
    }
}
\end{jcode}

\begin{concept}{Establish invariants in the constructor}
\index{invariant}
Why compute the subkeys in the constructor? Partly efficiency --- they are a pure
function of the key, and recomputing them per call would waste an AES operation
on every MAC.

But the deeper reason is \textbf{the invariant}. The constructor is the one place
that can guarantee ``this object is valid''. It checks the key length and derives
the subkeys; afterwards the object cannot be in a bad state, because nothing can
change it.

Contrast a design with a \code{setKey()} method. Now every method must ask: has
the key been set? Are the subkeys stale? Those questions do not exist if you make
them impossible.

\textbf{Establish invariants in the constructor, make fields \code{final},
expose no setters.} Half of ``good OOP'' is that one habit.
\end{concept}

\section{Verify --- RFC 4493 vectors}

These come from the RFC itself: independent of you, independent of this book,
and definitive.

\textbf{Key} $K = $ \hex{2b7e151628aed2a6abf7158809cf4f3c}

\begin{center}\small
\begin{tabular}{@{}L{0.44\linewidth}L{0.46\linewidth}@{}}
\toprule
\thead{Message} & \thead{Expected CMAC} \\
\midrule
\emph{(empty)} & \hex{bb1d6929e95937287fa37d129b756746} \\
\hex{6bc1bee22e409f96e93d7e117393172a} & \hex{070a16b46b4d4144f79bdd9dd04a287c} \\
\hex{6bc1bee22e409f96e93d7e117393172a}\newline
\hex{ae2d8a571e03ac9c9eb76fac45af8e51}\newline
\hex{30c81c46a35ce411} \emph{(40 B)} & \hex{dfa66747de9ae63030ca32611497c827} \\
the above $+$ \hex{e5fbc1191a0a52ef}\newline
\hex{f69f2445df4f9b17ad2b417be66c3710} \emph{(64 B)} & \hex{51f0bebf7e3b9d92fc49741779363cfe} \\
\bottomrule
\end{tabular}
\end{center}

Those four are chosen to hit every branch: empty (pad, $K_2$), exactly one block
($K_1$), a partial last block (pad, $K_2$), and multiple full blocks ($K_1$). If
all four pass, your implementation is right.

The subkeys, useful mid-debug:

\begin{trace}
K1 = fbeed618357133667c85e08f7236a8de
K2 = f7ddac306ae266ccf90bc11ee46d513b
\end{trace}

\begin{jcode}
package com.midnightbrewer.hardware.crypto;

import org.junit.jupiter.api.Test;
import java.util.HexFormat;
import static org.junit.jupiter.api.Assertions.assertEquals;

class AesCmacTest {

    private static final byte[] K =
            HexFormat.of().parseHex("2b7e151628aed2a6abf7158809cf4f3c");

    private static String mac(String messageHex) {
        byte[] m = HexFormat.of().parseHex(messageHex);
        return HexFormat.of().formatHex(new AesCmac(K).compute(m));
    }

    @Test
    void rfc4493_example1_empty() {
        assertEquals("bb1d6929e95937287fa37d129b756746", mac(""));
    }

    @Test
    void rfc4493_example2_oneBlock() {
        assertEquals("070a16b46b4d4144f79bdd9dd04a287c",
                mac("6bc1bee22e409f96e93d7e117393172a"));
    }

    @Test
    void rfc4493_example3_partialLastBlock() {
        assertEquals("dfa66747de9ae63030ca32611497c827",
                mac("6bc1bee22e409f96e93d7e117393172a"
                  + "ae2d8a571e03ac9c9eb76fac45af8e51"
                  + "30c81c46a35ce411"));
    }

    @Test
    void rfc4493_example4_fullBlocks() {
        assertEquals("51f0bebf7e3b9d92fc49741779363cfe",
                mac("6bc1bee22e409f96e93d7e117393172a"
                  + "ae2d8a571e03ac9c9eb76fac45af8e51"
                  + "30c81c46a35ce411e5fbc1191a0a52ef"
                  + "f69f2445df4f9b17ad2b417be66c3710"));
    }
}
\end{jcode}

\begin{verifybox}
\begin{shell}
mvn test -Dtest=AesCmacTest      # all four must be green
\end{shell}
Those four numbers were published in 2006 by people who have never seen your
code. If your code reproduces them exactly, that is not ``it seems to work'' ---
that is proof.
\end{verifybox}

\code{HexFormat} is a Java 17+ built-in, so no dependency is needed for hex
parsing and printing --- which you will do constantly from here on.

\section{Traps}

\begin{traptable}
Empty-message vector fails, others pass & Missing the \code{n=1, lastComplete=false} special case. \\
Partial-block vectors fail & Padding wrong: it is \hex{0x80} then zeros, not all zeros. \\
Full-block vectors fail & Using $K_2$ where $K_1$ belongs, or vice versa. \\
All four fail, all plausible-looking & \code{leftShift} carry direction --- you must walk backwards. \\
Sporadic wrong results & You mutated \code{key} or a subkey. They must be write-once. \\
\end{traptable}

\section{Commit}

\begin{shell}
git add -A
git commit -m "M7a: AES-CMAC (RFC 4493), verified against the RFC vectors"
\end{shell}
